\chapter{Types and Values}

\hayashi{} is dynamically typed: variables have no declared type, but each
value carries its type at runtime. The type system is consistent —
operations between incompatible types produce clear errors rather than
silent coercions.

\section{Primitive types}

\subsection{Int}

64-bit signed integer. Integer literals are written without a decimal point:

\begin{lstlisting}
let n        = 42
let negative = -7
let large    = 1_000_000    // _ as visual separator
\end{lstlisting}

\subsection{Float}

64-bit floating point (IEEE 754). Requires a decimal point or scientific
notation:

\begin{lstlisting}
let pi    = 3.14159
let e     = 2.71828
let small = 1.5e-10
\end{lstlisting}

Arithmetic between \texttt{Int} and \texttt{Float} promotes to
\texttt{Float}:

\begin{lstlisting}
let x = 3 / 2      // 1    (integer division)
let y = 3 / 2.0    // 1.5
let z = 3.0 / 2    // 1.5
\end{lstlisting}

Dividing a \texttt{Float} by zero returns \texttt{inf} or \texttt{-inf},
never an error — following IEEE 754:

\begin{lstlisting}
display 1.0 / 0.0    // inf
display -1.0 / 0.0   // -inf
display 0.0 / 0.0    // nan
\end{lstlisting}

\begin{nota}
  Dividing an \texttt{Int} by zero produces a runtime error. Dividing a
  \texttt{Float} by zero follows IEEE 754 and returns \texttt{inf}.
\end{nota}

\subsection{Bool}

True or false. Literals: \lstinline{true} and \lstinline{false}.

\begin{lstlisting}
let active = true
let closed = false
\end{lstlisting}

Logical operators: \lstinline{and}, \lstinline{or}, \lstinline{not}.

\subsection{Nil}

Absence of value. Some functions return \lstinline{nil} when there is no
relevant result (for example, \lstinline{push}, which mutates a list in
place).

\begin{lstlisting}
let nothing = nil
display nothing    // nil
\end{lstlisting}

\section{Strings}

Strings are Unicode character sequences delimited by double quotes:

\begin{lstlisting}
let name = "Hayashi"
let path = "/home/user/data.csv"
\end{lstlisting}

Escape sequences: \texttt{\textbackslash n} (newline),
\texttt{\textbackslash t} (tab),
\texttt{\textbackslash"} (quote),
\texttt{\textbackslash\textbackslash} (backslash).

F-strings allow expression interpolation:

\begin{lstlisting}
let year = 2026
display f"Published in {year}"   // Published in 2026
\end{lstlisting}

\section{Lists}

Ordered collections of values of any type, delimited by brackets:

\begin{lstlisting}
let numbers = [1, 2, 3, 4, 5]
let mixed   = [1, "two", true, nil]
let empty   = []
\end{lstlisting}

\section{Dictionaries}

Key-value mappings. Keys and values may be of any type:

\begin{lstlisting}
let config = {"host": "localhost", "port": 5432}
let empty  = {}
\end{lstlisting}

Access by key:

\begin{lstlisting}
config["host"]      // "localhost"
\end{lstlisting}

\section{DataFrame}

The central type of \hayashi{}. A DataFrame is a two-dimensional table
with named columns. Each column is a vector of \texttt{Float} (missing
values represented by \texttt{NaN}).

\begin{lstlisting}[caption={Creating a DataFrame inline}]
input x y
  1  2
  3  4
  5  6
end
\end{lstlisting}

DataFrames are discussed in depth in Chapter~\ref{cap:dataframes}.

\section{Conversions}

\hayashi{} does not perform implicit coercions between types. Conversions
must be explicit:

\begin{lstlisting}
let s = str(42)        // "42"
let n = int("7")       // 7
let f = float("3.14")  // 3.14
\end{lstlisting}

Attempting to add an \texttt{Int} to a \texttt{String} produces a type
error:

\begin{lstlisting}
42 + "3"
// error: type mismatch — expected Int or Float, got String
\end{lstlisting}

\section{Type Checking}

To dynamically check the type of a value, \hayashi{} provides the following built-in functions returning a boolean:

\begin{lstlisting}
is_nil(nil)         // true
is_int(42)          // true
is_float(3.14)      // true
is_bool(true)       // true
is_str("hello")     // true  (or is_string)
is_list([1, 2])     // true
is_dict({"a": 1})   // true
is_df(df)           // true  (or is_dataframe)
is_fn(|x| x)        // true  (or is_function)
\end{lstlisting}

Additionally, the \lstinline{type(x)} (or \lstinline{typeof(x)}) function returns the type of the value as a string (e.g. \lstinline{"int"}, \lstinline{"float"}, \lstinline{"nil"}, etc.).
